% GENERATED FILE. Edit canonical lessons, then run npm run generate:books.
% canonical-source-hash: fnv1a64:4658cd141ec871fb
% canonical-lessons: ML-C380-alarcca, ML-C380-kulappam, ML-C380-nilavili, ML-C380-nilavilikkuka, ML-C380-pintiriyuka, ML-R380-review-373-380

\chapter{Screams in the Night}
\label{ch:ml-screams-in-the-night}
\hlchaptermodality{380}

\begin{chapteropening}
\textbf{By the end of this chapter:} \emph{I can say a roar, trouble, a scream, to scream and to turn back.}

Complete the last lesson of chapter 380: I can say a roar, trouble, a scream, to scream and to turn back.
\end{chapteropening}

\section[\texorpdfstring{\ml{അലർച്ച} (alarcca)}{അലർച്ച (alarcca)}]{\ml{അലർച്ച} (alarcca) — a roar; a bellow}

\label{lesson:ML-C380-alarcca}



\begin{quote}
Before the new one: say the Malayalam for to overturn, then the Malayalam for to trip and fall.
\end{quote}

\subsection*{You'll want to know: \ml{അലർച്ച}}
\textbf{\ml{അലർച്ച}} — \emph{alarcca} — \textquotedblleft{}a roar; a bellow\textquotedblright{}.

\begin{quote}
\ml{രാത്രി} \ml{കാട്ടിൽ} \ml{നിന്ന്} \ml{ഒരു} \ml{കടുവയുടെ} \ml{അലർച്ച} \ml{കേട്ടു}.
\end{quote}

\begin{quote}
\emph{rātri kāṭṭil ninnŭ oru kaṭuvayuṭe alarcca kēṭṭu.}
\end{quote}

\begin{quote}
At night a tiger's roar was heard from the forest.
\end{quote}

The cry of a tiger, or of an angry crowd.

\begin{cousinweb}
\textbf{\ml{അലർച്ച}} (\emph{alarcca}) is native Dravidian, the noun of \ml{അലറുക}, which DEDR 245 lists as \textquotedblleft{}to roar, to bellow, to cry\textquotedblright{}.
\end{cousinweb}

\begin{grammarlens}[title={What you've built}]
One more word for telling what happened.
\end{grammarlens}

\subsection*{Your turn}
\begin{itemize}
\raggedright
  \item \emph{Say it:} \emph{alarcca}
  \item \emph{Say it:} \emph{alarcca}, once more
  \item \emph{Say it:} \emph{rātri kāṭṭil ninnŭ oru kaṭuvayuṭe alarcca kēṭṭu.} — At night a tiger's roar was heard from the forest.
  \item \emph{Recall:} say the Malayalam for to overturn, then the Malayalam for to trip and fall, then say \emph{alarcca} again
\end{itemize}

\begin{tcolorbox}[breakable,colback=teal!4,colframe=teal!35!black,title={Before you move on}]
What does \ml{അലർച്ച} mean? (\textquotedblleft{}A roar; a bellow\textquotedblright{}.) Say it once more.
\end{tcolorbox}
\section[\texorpdfstring{\ml{കുഴപ്പം} (kuḻappaṁ)}{കുഴപ്പം (kuḻappaṁ)}]{\ml{കുഴപ്പം} (kuḻappaṁ) — trouble; a problem}

\label{lesson:ML-C380-kulappam}



\begin{quote}
Before the new one: say the Malayalam for to trip and fall, then the Malayalam for a roar.
\end{quote}

\subsection*{You'll want to know: \ml{കുഴപ്പം}}
\textbf{\ml{കുഴപ്പം}} — \emph{kuḻappaṁ} — \textquotedblleft{}trouble; a problem\textquotedblright{}.

\begin{quote}
\ml{ഒന്നും} \ml{കുഴപ്പമില്ല}, \ml{എല്ലാവരും} \ml{സുരക്ഷിതരാണ്}.
\end{quote}

\begin{quote}
\emph{onnuṁ kuḻappamilla, ellāvaruṁ surakṣitarāṇŭ.}
\end{quote}

\begin{quote}
Nothing is wrong, everyone is safe.
\end{quote}

\ml{കുഴപ്പമില്ല} means no problem, it is all right.

\begin{cousinweb}
\textbf{\ml{കുഴപ്പം}} (\emph{kuḻappaṁ}) is native Dravidian: DEDR 1817 sets it beside Tamil \emph{kuḻappam}, \textquotedblleft{}confusion, agitation, disturbance\textquotedblright{}, from a verb meaning to mix up, to stir.
\end{cousinweb}

\begin{grammarlens}[title={What you've built}]
One more word for telling what happened.
\end{grammarlens}

\subsection*{Your turn}
\begin{itemize}
\raggedright
  \item \emph{Say it:} \emph{kuḻappaṁ}
  \item \emph{Say it:} \emph{kuḻappaṁ}, once more
  \item \emph{Say it:} \emph{onnuṁ kuḻappamilla, ellāvaruṁ surakṣitarāṇŭ.} — Nothing is wrong, everyone is safe.
  \item \emph{Recall:} say the Malayalam for to trip and fall, then the Malayalam for a roar, then say \emph{kuḻappaṁ} again
\end{itemize}

\begin{tcolorbox}[breakable,colback=teal!4,colframe=teal!35!black,title={Before you move on}]
What does \ml{കുഴപ്പം} mean? (\textquotedblleft{}Trouble; a problem\textquotedblright{}.) Say it once more.
\end{tcolorbox}
\section[\texorpdfstring{\ml{നിലവിളി} (nilaviḷi)}{നിലവിളി (nilaviḷi)}]{\ml{നിലവിളി} (nilaviḷi) — a scream; a cry for help}

\label{lesson:ML-C380-nilavili}



\begin{quote}
Before the new one: say the Malayalam for a roar, then the Malayalam for trouble.
\end{quote}

\subsection*{You'll want to know: \ml{നിലവിളി}}
\textbf{\ml{നിലവിളി}} — \emph{nilaviḷi} — \textquotedblleft{}a scream; a cry for help\textquotedblright{}.

\begin{quote}
\ml{വീട്ടിൽ} \ml{നിന്ന്} \ml{ഒരു} \ml{നിലവിളി} \ml{കേട്ടു}.
\end{quote}

\begin{quote}
\emph{vīṭṭil ninnŭ oru nilaviḷi kēṭṭu.}
\end{quote}

\begin{quote}
A scream was heard from the house.
\end{quote}

Often a cry of grief or of fear.

\begin{cousinweb}
\textbf{\ml{നിലവിളി}} (\emph{nilaviḷi}) is native Dravidian, usually traced to \ml{നില}, \textquotedblleft{}a standing, a place\textquotedblright{} (DEDR 3675), and \ml{വിളി} (\emph{viḷi}), \textquotedblleft{}a call, a cry\textquotedblright{} (DEDR 5433).
\end{cousinweb}

\begin{grammarlens}[title={What you've built}]
One more word for telling what happened.
\end{grammarlens}

\subsection*{Your turn}
\begin{itemize}
\raggedright
  \item \emph{Say it:} \emph{nilaviḷi}
  \item \emph{Say it:} \emph{nilaviḷi}, once more
  \item \emph{Say it:} \emph{vīṭṭil ninnŭ oru nilaviḷi kēṭṭu.} — A scream was heard from the house.
  \item \emph{Recall:} say the Malayalam for a roar, then the Malayalam for trouble, then say \emph{nilaviḷi} again
\end{itemize}

\begin{tcolorbox}[breakable,colback=teal!4,colframe=teal!35!black,title={Before you move on}]
What does \ml{നിലവിളി} mean? (\textquotedblleft{}A scream; a cry for help\textquotedblright{}.) Say it once more.
\end{tcolorbox}
\section[\texorpdfstring{\ml{നിലവിളിക്കുക} (nilaviḷikkuka)}{നിലവിളിക്കുക (nilaviḷikkuka)}]{\ml{നിലവിളിക്കുക} (nilaviḷikkuka) — to scream; to cry out}

\label{lesson:ML-C380-nilavilikkuka}



\begin{quote}
Before the new one: say the Malayalam for trouble, then the Malayalam for a scream.
\end{quote}

\subsection*{You'll want to know: \ml{നിലവിളിക്കുക}}
\textbf{\ml{നിലവിളിക്കുക}} — \emph{nilaviḷikkuka} — \textquotedblleft{}to scream; to cry out\textquotedblright{}.

\begin{quote}
\ml{പാമ്പിനെ} \ml{കണ്ട്} \ml{കുട്ടി} \ml{നിലവിളിച്ചു}.
\end{quote}

\begin{quote}
\emph{pāmpine kaṇṭŭ kuṭṭi nilaviḷiccu.}
\end{quote}

\begin{quote}
Seeing the snake, the child screamed.
\end{quote}

Past: \ml{നിലവിളിച്ചു}, \emph{nilaviḷiccu}.

Louder than \ml{കരയുക}, to cry: a cry for help.

\begin{cousinweb}
\textbf{\ml{നിലവിളിക്കുക}} (\emph{nilaviḷikkuka}) is native Dravidian: the verb of \ml{നിലവിളി} (\emph{nilaviḷi}), \textquotedblleft{}a scream\textquotedblright{}, ending in \ml{വിളിക്കുക}, \textquotedblleft{}to call\textquotedblright{} (DEDR 5433).
\end{cousinweb}

\begin{grammarlens}[title={What you've built}]
One more word for telling what happened.
\end{grammarlens}

\subsection*{Your turn}
\begin{itemize}
\raggedright
  \item \emph{Say it:} \emph{nilaviḷikkuka}
  \item \emph{Say it:} \emph{nilaviḷikkuka}, once more
  \item \emph{Say it:} \emph{pāmpine kaṇṭŭ kuṭṭi nilaviḷiccu.} — Seeing the snake, the child screamed.
  \item \emph{Recall:} say the Malayalam for trouble, then the Malayalam for a scream, then say \emph{nilaviḷikkuka} again
\end{itemize}

\begin{tcolorbox}[breakable,colback=teal!4,colframe=teal!35!black,title={Before you move on}]
What does \ml{നിലവിളിക്കുക} mean? (\textquotedblleft{}To scream; to cry out\textquotedblright{}.) Say it once more.
\end{tcolorbox}
\section[\texorpdfstring{\ml{പിന്തിരിയുക} (pintiriyuka)}{പിന്തിരിയുക (pintiriyuka)}]{\ml{പിന്തിരിയുക} (pintiriyuka) — to turn back; to retreat}

\label{lesson:ML-C380-pintiriyuka}



\begin{quote}
Before the new one: say the Malayalam for a scream, then the Malayalam for to scream.
\end{quote}

\subsection*{You'll want to know: \ml{പിന്തിരിയുക}}
\textbf{\ml{പിന്തിരിയുക}} — \emph{pintiriyuka} — \textquotedblleft{}to turn back; to retreat\textquotedblright{}.

\begin{quote}
\ml{ആനയെ} \ml{കണ്ട്} \ml{ഞങ്ങൾ} \ml{പിന്തിരിഞ്ഞു}.
\end{quote}

\begin{quote}
\emph{ānaye kaṇṭŭ ñaṅṅaḷ pintiriññu.}
\end{quote}

\begin{quote}
Seeing the elephant, we turned back.
\end{quote}

Past: \ml{പിന്തിരിഞ്ഞു}, \emph{pintiriññu}.

\ml{പിൻ}, back, + \ml{തിരിയുക}, to turn.

\begin{cousinweb}
\textbf{\ml{പിന്തിരിയുക}} (\emph{pintiriyuka}) is a native Dravidian compound of \ml{പിൻ} (\emph{pin}), \textquotedblleft{}back, behind\textquotedblright{} (DEDR 4205), and \ml{തിരിയുക} (\emph{tiriyuka}), \textquotedblleft{}to turn\textquotedblright{} (DEDR 3246).
\end{cousinweb}

\begin{grammarlens}[title={What you've built}]
One more word for telling what happened.
\end{grammarlens}

\subsection*{Your turn}
\begin{itemize}
\raggedright
  \item \emph{Say it:} \emph{pintiriyuka}
  \item \emph{Say it:} \emph{pintiriyuka}, once more
  \item \emph{Say it:} \emph{ānaye kaṇṭŭ ñaṅṅaḷ pintiriññu.} — Seeing the elephant, we turned back.
  \item \emph{Recall:} say the Malayalam for a scream, then the Malayalam for to scream, then say \emph{pintiriyuka} again
\end{itemize}

\begin{tcolorbox}[breakable,colback=teal!4,colframe=teal!35!black,title={Before you move on}]
What does \ml{പിന്തിരിയുക} mean? (\textquotedblleft{}To turn back; to retreat\textquotedblright{}.) Say it once more.
\end{tcolorbox}
\section[(dialogue)]{Words from chapters 373–380 — a review}

\label{lesson:ML-R380-review-373-380}



\begin{quote}
Say the Malayalam for an attack and a threat.
\end{quote}

\subsection*{Your turn}
Say these aloud:

\begin{itemize}
\raggedright
  \item an attack, a threat, to attack, to loot, to smash
  \item a warning, a disaster, a rescue operation, a state of emergency, to evacuate
  \item a rescuer, a volunteer, to rescue, to escape, to get stuck
  \item a victim, life, a dead body, to get hurt, to drown
  \item an earthquake, a landslide, to be washed away, to be cut off, to go missing
  \item shelter, a refugee, loss of life, to flee, to seek
  \item a throng, to crash into, to collide, to overturn, to trip and fall
  \item a roar, trouble, a scream, to scream, to turn back
\end{itemize}

\begin{tcolorbox}[breakable,colback=teal!4,colframe=teal!35!black,title={Before you move on}]
An attack? (\textbf{\ml{ആക്രമണം}}, \emph{ākramaṇaṁ}.) A threat? (\textbf{\ml{ഭീഷണി}}, \emph{bhīṣaṇi}.) To attack? (\textbf{\ml{ആക്രമിക്കുക}}, \emph{ākramikkuka}.) To loot? (\textbf{\ml{കൊള്ളയടിക്കുക}}, \emph{koḷḷayaṭikkuka}.) To smash? (\textbf{\ml{തകർക്കുക}}, \emph{takarkkuka}.) A warning? (\textbf{\ml{മുന്നറിയിപ്പ്}}, \emph{munnaṟiyippŭ}.) A disaster? (\textbf{\ml{ദുരന്തം}}, \emph{durantaṁ}.) A rescue operation? (\textbf{\ml{രക്ഷാപ്രവർത്തനം}}, \emph{rakṣāpravarttanaṁ}.) A state of emergency? (\textbf{\ml{അടിയന്തരാവസ്ഥ}}, \emph{aṭiyantarāvastha}.) To evacuate? (\textbf{\ml{ഒഴിപ്പിക്കുക}}, \emph{oḻippikkuka}.) A rescuer? (\textbf{\ml{രക്ഷകൻ}}, \emph{rakṣakan}.) A volunteer? (\textbf{\ml{സന്നദ്ധപ്രവർത്തകൻ}}, \emph{sannaddhapravarttakan}.) To rescue? (\textbf{\ml{രക്ഷിക്കുക}}, \emph{rakṣikkuka}.) To escape? (\textbf{\ml{രക്ഷപ്പെടുക}}, \emph{rakṣappeṭuka}.) To get stuck? (\textbf{\ml{കുടുങ്ങുക}}, \emph{kuṭuṅṅuka}.) A victim? (\textbf{\ml{ഇര}}, \emph{ira}.) Life? (\textbf{\ml{ജീവൻ}}, \emph{jīvan}.) A dead body? (\textbf{\ml{മൃതദേഹം}}, \emph{mṛtadēhaṁ}.) To get hurt? (\textbf{\ml{പരിക്കേൽക്കുക}}, \emph{parikkēlkkuka}.) To drown? (\textbf{\ml{മുങ്ങിമരിക്കുക}}, \emph{muṅṅimarikkuka}.) An earthquake? (\textbf{\ml{ഭൂകമ്പം}}, \emph{bhūkampaṁ}.) A landslide? (\textbf{\ml{ഉരുൾപൊട്ടൽ}}, \emph{uruḷpoṭṭal}.) To be washed away? (\textbf{\ml{ഒലിച്ചുപോകുക}}, \emph{oliccupōkuka}.) To be cut off? (\textbf{\ml{ഒറ്റപ്പെടുക}}, \emph{oṟṟappeṭuka}.) To go missing? (\textbf{\ml{കാണാതാകുക}}, \emph{kāṇātākuka}.) Shelter? (\textbf{\ml{അഭയം}}, \emph{abhayaṁ}.) A refugee? (\textbf{\ml{അഭയാർത്ഥി}}, \emph{abhayārtthi}.) Loss of life? (\textbf{\ml{ജീവഹാനി}}, \emph{jīvahāni}.) To flee? (\textbf{\ml{പലായനം} \ml{ചെയ്യുക}}, \emph{palāyanaṁ ceyyuka}.) To seek? (\textbf{\ml{തേടുക}}, \emph{tēṭuka}.) A throng? (\textbf{\ml{ആൾക്കൂട്ടം}}, \emph{āḷkkūṭṭaṁ}.) To crash into? (\textbf{\ml{ഇടിക്കുക}}, \emph{iṭikkuka}.) To collide? (\textbf{\ml{കൂട്ടിമുട്ടുക}}, \emph{kūṭṭimuṭṭuka}.) To overturn? (\textbf{\ml{മറിയുക}}, \emph{maṟiyuka}.) To trip and fall? (\textbf{\ml{തട്ടിവീഴുക}}, \emph{taṭṭivīḻuka}.) A roar? (\textbf{\ml{അലർച്ച}}, \emph{alarcca}.) Trouble? (\textbf{\ml{കുഴപ്പം}}, \emph{kuḻappaṁ}.) A scream? (\textbf{\ml{നിലവിളി}}, \emph{nilaviḷi}.) To scream? (\textbf{\ml{നിലവിളിക്കുക}}, \emph{nilaviḷikkuka}.) To turn back? (\textbf{\ml{പിന്തിരിയുക}}, \emph{pintiriyuka}.)
\end{tcolorbox}
